% =====================================================================
%  paper.tex  —  C2 挑战交付物
%  编译方式：xelatex -> bibtex -> xelatex -> xelatex
%  （或直接使用 latexmk -xelatex paper.tex / tectonic paper.tex）
% =====================================================================
\documentclass[UTF8,a4paper,11pt]{ctexart}

% ---------- 基础排版 ----------
\usepackage[margin=2.5cm]{geometry}
\usepackage{amsmath,amssymb,amsthm}
\usepackage{graphicx}
\usepackage{booktabs}
\usepackage{array}
\usepackage{multirow}
\usepackage{enumitem}
\usepackage{float}
\usepackage{caption}
\usepackage{xcolor}
\usepackage{tikz}
\usetikzlibrary{arrows.meta,positioning,shapes.geometric,fit,backgrounds}
\usepackage[hidelinks]{hyperref}
\usepackage[numbers,sort&compress]{natbib}
\usepackage{algorithm}
\usepackage{algpseudocode}

\captionsetup{font=small,labelfont=bf}
\bibliographystyle{plainnat}

% ---------- 定理环境 ----------
\newtheorem{definition}{定义}[section]
\newtheorem{proposition}{命题}[section]
\newtheorem{remark}{注记}[section]

% ---------- 自定义命令 ----------
% 语义括号（不依赖 stmaryrd，避免额外宏包）
\providecommand{\llbracket}{\mathopen{[\![}}
\providecommand{\rrbracket}{\mathclose{]\!]}}
\newcommand{\NL}{\mathrm{NL}}
\newcommand{\IR}{\mathrm{IR}}
\newcommand{\FOL}{\mathrm{FOL}}
\newcommand{\HOL}{\mathrm{HOL}}
\newcommand{\sem}[1]{\llbracket #1 \rrbracket}
\newcommand{\code}[1]{\texttt{#1}}
\DeclareMathOperator{\Verify}{Verify}
\DeclareMathOperator{\Reduce}{Reduce}
\DeclareMathOperator{\Render}{Render}

% ---------- 标题信息 ----------
\title{\bfseries 从自然语言到可验证推理：\\ 面向可靠 AI 数学的语义规约框架}

\author{%
  作者姓名\thanks{本文为 AI 挑战任务 C2 交付物。写作大纲依据课程一手材料
  《Lectures on AI for Mathematics》（同济大学，2026）\citep{chen2026lectures}，
  全部引用文献均经联网核实，见 \code{references.bib}。}%
  \\[2pt]
  \small 单位名称 \quad \texttt{author@example.edu}
}

\date{\today}

\begin{document}
\maketitle

% =====================================================================
\begin{abstract}
\noindent
大语言模型（LLM）在数学推理任务上的表现持续提升，但其输出以自然语言为载体，
既不可直接验证，也无法稳定传导到形式化证明系统。本文提出一个统一的
\textbf{三层分解模型}，把「可靠的 AI 数学推理」拆解为\emph{生成层}
（Generation Layer）、\emph{语义规约层}（Semantic Reduction Layer）与
\emph{验证层}（Verification Layer）。基于对现有系统的分析，我们论证：
当前系统的主要失败来源并非生成能力不足，也非验证器不可靠，而是
\emph{语义规约层}的缺失与不稳定——自然语言命题到形式对象之间的映射既
不具备单值性（ambiguity），也不具备类型完整性（missing typing），更不具备
可逆性（non-invertibility）。据此，本文提出以\textbf{类型化中间表示}
（Typed Intermediate Representation, Typed IR）为核心的框架：将自然语言
仅作为人机接口，把可验证结构作为执行内核，并通过\textbf{双向一致性}
（Bidirectional Consistency）约束保证 IR 与自然语言之间的可解释性与可回溯性。
我们给出框架的形式化定义、流水线架构（生成 $\to$ IR $\to$ 验证 $\to$ 渲染）
以及可证伪的实验设计，并讨论其在自动定理证明、数学解题与 AI 教育三类场景中的
应用。本文的贡献是概念性的：它把「语义规约」从工程细节提升为核心系统组件，
并给出与之匹配的评估指标（可验证率、一致性、回译保真度）。

\vspace{6pt}
\noindent\textbf{关键词：} AI 数学；语义规约；类型化中间表示；自动形式化；
神经符号系统；Curry--Howard 对应；形式验证
\end{abstract}

% =====================================================================
\section{引言}\label{sec:intro}

\subsection{背景与动机}

近年来的经验证据表明，大语言模型已经具备可观的数学推理能力。在
Chain-of-Thought（CoT）提示下，模型能够生成人类可读的多步推理链
\citep{wei2022cot}；在 GSM8K \citep{cobbe2021gsm8k} 与 MATH
\citep{hendrycks2021math} 等基准上，准确率随模型规模与推理时计算量的增加
而显著提高。与此同时，将计算外包给程序解释器（Program-of-Thoughts, PoT）
\citep{chen2023pot}、或让模型自行学习调用外部工具（Toolformer）
\citep{schick2023toolformer}，也都带来了可测量的收益。

然而，这些能力的「可靠性」是脆弱的。自然语言推理链是一个
\emph{非类型化、概率性}的文本序列：它的每一步都不携带形式语义，因而
\begin{enumerate}[leftmargin=1.6em,itemsep=2pt]
  \item \textbf{不可直接验证}——只存在「答案对不对」这一末端信号，
        中间步骤的正确性无法被机械地检查；
  \item \textbf{不可局部定位}——当结论错误时，无法判定失败发生在哪一步，
        因为步骤之间没有可推导关系；
  \item \textbf{不可迁移}——同一条推理链无法被复用到形式系统中，
        也无法作为证明对象被归档、检索与组合。
\end{enumerate}

与之相对，形式数学系统（Lean 4 \citep{demoura2021lean4}、Coq、Isabelle/HOL、
SMT 求解器如 Z3 \citep{demoura2008z3}）提供了\emph{机器可检查}的证明语义：
一个被编译器接受的证明对象，其正确性不依赖于生成它的模型是否「说谎」。
AlphaGeometry \citep{trinh2024alphageometry}、FunSearch
\citep{romera2024funsearch} 与 AlphaProof \citep{hubert2025alphaproof}
正是这一路线的代表性成果：AlphaProof 在 2024 年国际数学奥林匹克（IMO）
上结合 AlphaGeometry 解出六题中的四题，达到银牌水平，且其推理被形式系统自动
验证。

\subsection{问题陈述}

于是问题自然浮现：\emph{既然两端都足够强，为什么端到端的可靠 AI 数学仍然困难？}

本文的答案是：\textbf{缺失的是中间那一段}。生成端（LLM）与验证端
（形式系统）之间存在一个语义鸿沟——自然语言命题需要被翻译成形式对象
（这一过程称为\emph{自动形式化}，autoformalization \citep{wu2022autoformalization}）。
这一翻译\emph{不是}一个轻量的预处理步骤，而是一个独立的、承载全部正确性风险的
系统组件。我们把这一组件称为\textbf{语义规约层}。

\subsection{核心论点与贡献}

本文的核心论点是：

\begin{quote}
\itshape
可靠 AI 数学推理 $\neq$ 生成 $+$ 验证。\\
可靠性不能仅依赖「最终验证」，必须显式引入\emph{语义规约层}作为中间锚点。
\end{quote}

具体贡献如下：
\begin{enumerate}[leftmargin=1.6em,itemsep=2pt]
  \item \textbf{概念贡献}：提出三层分解模型（\S\ref{sec:model}），
        把「语义规约」确立为独立且核心的系统层，并论证它是当前系统的主要失败来源
        （\S\ref{sec:bottleneck}）。
  \item \textbf{方法贡献}：给出以 Typed IR 为核心的框架设计
        （\S\ref{sec:framework}），包括类型系统、操作语义、约束系统与
        可组合结构，以及「生成 $\to$ IR $\to$ 验证 $\to$ 渲染」的标准流水线。
  \item \textbf{理论视角}：以 Curry--Howard 对应为语义规约提供类型论基础，
        并形式化定义\emph{双向一致性}约束（可解释性 + 可回溯性）。
  \item \textbf{评估贡献}：提出可证伪的实验设计（\S\ref{sec:experiment}），
        引入比 accuracy 更敏感的指标——\emph{可验证率}与\emph{回译保真度}。
\end{enumerate}

本文是\emph{立场与方法}论文（position-and-method paper）：我们提出的框架
以设计论证与可检验假设的形式给出，实验部分为预注册式的设计说明，
而非已完成的实证结果。\S\ref{sec:discussion} 明确讨论了这一局限。

% =====================================================================
\section{问题分层模型}\label{sec:model}

\subsection{三层结构}

我们把一个 AI 数学推理系统分解为三个层次，如表~\ref{tab:layers} 所示。

\begin{table}[htbp]
\centering
\caption{三层分解模型：性质、契约与代表工具}
\label{tab:layers}
\small
\begin{tabular}{@{}p{2.3cm}p{3.4cm}p{3.6cm}p{4.2cm}@{}}
\toprule
\textbf{层} & \textbf{输出} & \textbf{性质} & \textbf{代表工具/方法} \\
\midrule
生成层
  & 自然语言推理链
  & 概率性、非类型化、非严格语义
  & CoT \citep{wei2022cot}；PoT \citep{chen2023pot}；Toolformer \citep{schick2023toolformer} \\
\addlinespace
语义规约层
  & 结构化形式对象
  & \textbf{承载正确性风险的关键层}
  & 自动形式化 \citep{wu2022autoformalization}；Typed IR（本文）\\
\addlinespace
验证层
  & 布尔判定 / 证明对象
  & 可靠、可判定（在系统内）、不可欺骗
  & Lean 4 \citep{demoura2021lean4}；Z3 \citep{demoura2008z3}；Coq；Agda \\
\bottomrule
\end{tabular}
\end{table}

\subsection{第一层：生成层（Generation Layer）}

生成层接收自然语言问题 $q \in \mathcal{Q}_{\NL}$，产出推理文本
$r = (s_1, s_2, \dots, s_n)$，其中每个 $s_i$ 是一段自然语言陈述。其性质为：
\begin{itemize}[leftmargin=1.6em,itemsep=2pt]
  \item \textbf{概率性}：$r \sim P_\theta(\cdot \mid q)$，同一输入可产生不同输出；
  \item \textbf{非类型化}：不存在 $\vdash s_i : T$ 这样的类型判定关系；
  \item \textbf{非严格语义}：$s_i$ 的真值不独立于上下文与世界知识。
\end{itemize}

CoT \citep{wei2022cot} 展示了该层的核心价值：显式中间步骤可以显著提升
多步推理的准确性。但这同时暴露了该层的根本缺陷——\emph{步骤的「看起来合理」
与「语义正确」之间没有保障}。PoT \citep{chen2023pot} 与 Toolformer
\citep{schick2023toolformer} 部分缓解了这一点：前者把\emph{计算}从
\emph{推理}中剥离，交给外部解释器执行；后者让模型学习在何时调用哪个工具。
二者的共同思路是「把不可靠的环节外包给可靠系统」——这正是本文所称
语义规约的雏形，但它们都只覆盖了\emph{计算}这一个侧面，未触及
\emph{语义}。

\subsection{第二层：语义规约层（Semantic Reduction Layer）}

语义规约层的职责是把自然语言陈述映射为形式对象，即定义一个映射
\begin{equation}\label{eq:reduce}
  \Reduce : \mathcal{Q}_{\NL} \times \mathcal{R}_{\NL}
            \;\longrightarrow\; \mathcal{T}_{\IR},
\end{equation}
其中 $\mathcal{R}_{\NL}$ 是推理文本空间，$\mathcal{T}_{\IR}$ 是类型化中间
表示的集合。$\mathcal{T}_{\IR}$ 中的元素可以是
\begin{itemize}[leftmargin=1.6em,itemsep=2pt]
  \item 逻辑表达式（FOL / HOL 公式）；
  \item 可执行程序（Python / 领域特定语言 DSL）；
  \item 证明对象（proof objects）；
  \item 类型化中间表示（Typed IR），即本文主张的统一载体。
\end{itemize}

该层的理论基础是 \textbf{Curry--Howard 对应}（命题即类型、证明即程序）：
它保证了「构造一个类型为 $T$ 的项」与「证明命题 $T$」是同构的。因此，一旦
自然语言命题被规约为类型 $T$，验证问题就化为类型检查问题，而验证层的
可靠性可以完全继承。

但严格的对应关系也立刻暴露出该层的困难：自然语言是\emph{欠类型化}的，
而类型系统要求\emph{完全类型化}——这个落差必须在该层内部被弥合，
且无法由下游的验证层代劳。这正是 \S\ref{sec:bottleneck} 要展开的核心。

\subsection{第三层：验证层（Verification Layer）}

验证层接收形式对象，返回判定结果。其检查内容包括：
\begin{itemize}[leftmargin=1.6em,itemsep=2pt]
  \item \textbf{自洽性}（consistency）：对象内部无矛盾；
  \item \textbf{可推导性}（derivability）：证明对象可由公理与已有定理推出；
  \item \textbf{事实一致性}（factual grounding）：结论与已知数学事实相容；
  \item \textbf{约束满足}（constraint satisfaction）：所有显式约束被满足。
\end{itemize}

该层的工程可靠性已被充分验证：Lean 4 的 kernel
\citep{demoura2021lean4} 与 Z3 \citep{demoura2008z3} 都是成熟且久经考验的
系统；AlphaProof 的全部成绩都建立在「形式系统自动验证」之上
\citep{hubert2025alphaproof}。因此，\emph{验证层不是瓶颈}——这一判断对
后续的归因至关重要。

\begin{figure}[htbp]
\centering
\begin{tikzpicture}[
  node distance=8mm and 12mm,
  box/.style={draw, rounded corners=3pt, minimum height=1.0cm,
              minimum width=4.6cm, align=center, font=\small},
  layerbox/.style={draw, dashed, rounded corners=6pt, inner sep=6pt},
  ar/.style={-{Latex[length=2.5mm]}, thick}
]
\node[box, fill=blue!8]  (gen)  {\textbf{生成层}\\[1pt]
  \footnotesize 自然语言推理链 $r$\\ \footnotesize（概率性、非类型化）};
\node[box, fill=orange!12, below=of gen] (red) {\textbf{语义规约层}\\[1pt]
  \footnotesize Typed IR $\;\tau \in \mathcal{T}_{\IR}$\\ \footnotesize（瓶颈所在）};
\node[box, fill=green!10, below=of red]  (ver) {\textbf{验证层}\\[1pt]
  \footnotesize 类型检查 / SMT 判定\\ \footnotesize（可靠、可判定）};

\draw[ar] (gen) -- node[right, font=\footnotesize, xshift=2pt]{$\Reduce$（自动形式化）} (red);
\draw[ar] (red) -- node[right, font=\footnotesize, xshift=2pt]{$\Verify$（类型检查）} (ver);
\draw[ar, dashed, gray] (ver.east) to[out=0,in=0, looseness=1.9]
  node[right, font=\footnotesize, align=left, xshift=2pt]
  {失败回传\\（反例 $\to$ 修订 IR）} (red.east);

\node[layerbox, fit=(gen), inner sep=4pt] {};
\end{tikzpicture}
\caption{三层分解与语义规约层的中心地位。实线为前向流水线，
虚线为失败回传路径：验证失败时，修订的对象应是 IR 而非自然语言推理链——
这正是把语义规约确立为独立层的直接工程后果。}
\label{fig:three-layers}
\end{figure}

% =====================================================================
\section{核心问题：语义规约瓶颈}\label{sec:bottleneck}

\subsection{归因观察}

我们对现有系统的失败模式给出如下归因假设：

\begin{quote}
\itshape
当前系统的失败主要源于第二层（语义规约），而非生成层或验证层。
\end{quote}

这一归因可以被具体的经验事实支持。autoformalization 的原始工作报告：
LLM 能够把约 $25.3\%$ 的数学竞赛题\emph{完美}翻译为 Isabelle/HOL 形式规格
\citep{wu2022autoformalization}——换言之，\emph{近四分之三的翻译存在缺陷}。
与此同时，同一批工作显示，用自动形式化产物训练神经定理证明器后，
miniF2F \citep{zheng2022minif2f} 上的证明率从 $29.6\%$ 提升到 $35.2\%$。
这说明：\emph{涨点来自规约质量，而非生成模型容量}。

更近期的证据来自 PutnamBench \citep{tsoukalas2024putnambench}：这是一个
基于 Lean 4 的形式化奥林匹克级基准，其结果显示，即便在强模型上，
形式化竞赛数学的通过率仍然很低——失败更多发生在「把题面正确形式化」
这一环节，而非「找不到证明」。

\subsection{三点结构性困难}

我们把该层的困难归纳为三条：

\paragraph{(1) 语义不稳定（ambiguity）.}
自然语言命题与形式命题之间不存在确定映射。例如「$f$ 是连续的」
在 $\mathbb{R}$ 上至少对应 $\varepsilon$-$\delta$ 定义、序列连续、
拓扑连续三种形式化，而它们在一般拓扑空间中并不等价。$\Reduce$ 因此
\emph{不是函数}，而是一个关系 $\Reduce \subseteq \mathcal{Q}_{\NL} \times \mathcal{T}_{\IR}$。

\paragraph{(2) 类型不完整（missing typing）.}
自然语言不携带类型信息。量词的作用域、函数的定义域与值域、
集合的宇宙层级，均需由规约层补齐。若补齐错误，得到的形式对象
\emph{类型正确却语义错误}——这类错误最危险，因为它能通过验证层的
\emph{形式}检查，却验证了一个并非原问题的命题。

\paragraph{(3) 映射不可逆（non-invertibility）.}
$\Reduce$ 一般没有逆映射。这意味着无法仅凭 IR 恢复原始意图；
若缺少双向约束，系统将丧失可解释性与可回溯性——当验证通过时，
人类无法确认「通过的」究竟是原命题还是被悄悄改写过的命题。

\subsection{后果}

上述三点分别导致：
\begin{enumerate}[leftmargin=1.6em,itemsep=2pt]
  \item \textbf{验证失效}（no handle to verify）：$\Reduce$ 不确定时，
        验证层没有可检查的对象，只能退化为对最终答案的比对；
  \item \textbf{幻觉不可检测}：一个类型正确但语义漂移的对象会被验证器
        \emph{接受}，从而把幻觉「洗白」为已证结论；
  \item \textbf{错误不可归因}：由于缺少中间锚点，失败无法被定位到具体步骤，
        迭代改进只能靠盲试。
\end{enumerate}

\begin{remark}
第 (2) 点值得强调：\emph{验证器的可靠性与系统的可靠性不是同一件事}。
一个 100\% 正确的类型检查器，与一个把错误命题形式化的映射组合起来，
整体系统的可靠性可以接近零。这是「生成 $+$ 验证」这一朴素配方失效的
根本原因。
\end{remark}

% =====================================================================
\section{方法框架}\label{sec:framework}

\subsection{类型化中间表示（Typed IR）}

我们主张以统一的 Typed IR 作为语义规约层的输出载体。Typed IR 由四个
组件构成：

\paragraph{(a) 类型系统（Type System）.}
给定基类型集合 $\mathcal{B}$（如 $\mathbb{N}$、$\mathbb{R}$、$\mathrm{Prop}$）
与类型构造子（$\to$、$\times$、$\Pi$、$\Sigma$），定义类型项
\begin{equation}\label{eq:type}
  T \;::=\; b \mid T \to T \mid T \times T \mid \Pi x{:}T.\,T \mid \Sigma x{:}T.\,T,
  \qquad b \in \mathcal{B}.
\end{equation}

\paragraph{(b) 操作语义（Operational Semantics）.}
为 IR 的每个构造子给出小步语义规则 $\tau \longrightarrow \tau'$，使 IR
可以被\emph{执行}（用于程序类对象）或被\emph{归约}（用于证明类对象）。
这使得 IR 不是静态标注，而是可运行的语义载体。

\paragraph{(c) 约束系统（Constraints）.}
IR 上附加一组约束 $\mathcal{C} = \{c_1, \dots, c_m\}$，用于编码自然语言中
隐含但未言明的前提，例如定义域非空、量词作用域、可测性等。约束可被
翻译为 SMT 公式并交由 Z3 \citep{demoura2008z3} 判定。

\paragraph{(d) 可组合结构（Composable Graph / AST）.}
IR 是一个有向无环图（AST 或 workflow graph），节点是带类型的子对象，
边是依赖或蕴含关系。可组合性使得局部修订成为可能：验证失败时，
只需回退到失败的子节点重新规约，而不必重写整条推理链。

\begin{definition}[良类型 IR]\label{def:welltyped}
称 IR 对象 $\tau$ 关于上下文 $\Gamma$ 是\emph{良类型的}，记作
$\Gamma \vdash \tau : T$，当且仅当 $\tau$ 可由式~\eqref{eq:type} 的类型规则
推导，且其全部子节点良类型，且 $\mathcal{C}$ 在 $\tau$ 上可满足。
\end{definition}

\subsection{流水线架构}

框架的标准流水线为
\begin{equation}\label{eq:pipeline}
  \underbrace{q}_{\text{自然语言}}
  \;\xrightarrow{\;\Reduce\;}\;
  \underbrace{\tau}_{\text{Typed IR}}
  \;\xrightarrow{\;\Verify\;}\;
  \underbrace{v \in \{\top, \bot\}}_{\text{判定}}
  \;\xrightarrow{\;\Render\;}\;
  \underbrace{\hat{r}}_{\text{自然语言输出}} .
\end{equation}

关键设计原则是：

\begin{quote}
\itshape
自然语言只是 interface；可验证结构才是 execution core。
\end{quote}

即：人类读写自然语言，但系统内部的一切计算、传递与判定都发生在 IR 上。
$\Render$ 负责把验证后的 IR 重新渲染为人类可读的文本/证明说明，
但\emph{渲染结果不参与正确性判定}——它只是 IR 的一个视图。

\begin{algorithm}[htbp]
\caption{语义规约驱动的可靠推理（SRR: Semantic-Reduction Reasoning）}
\label{alg:srr}
\begin{algorithmic}[1]
\Require 自然语言问题 $q$；规约器 $\Reduce$；验证器 $\Verify$；渲染器 $\Render$
\Ensure 判定 $v$ 与解释 $\hat{r}$
\State $\tau_0 \gets \Reduce(q)$ \Comment{初次规约：生成候选 Typed IR}
\If{$\tau_0$ 未通过良类型检查（定义~\ref{def:welltyped}）}
    \State $\tau_0 \gets \textsc{Repair}(\tau_0)$ \Comment{补齐类型/约束，局部修订}
\EndIf
\State $\mathcal{C} \gets \textsc{ExtractConstraints}(\tau_0)$
\ForAll{$c \in \mathcal{C}$} \Comment{逐条 SMT 判定，定位失败约束}
    \If{$\textsc{SMT}(c) = \bot$}
        \State 记录反例 $\rho_c$；$\tau_0 \gets \textsc{Refine}(\tau_0, \rho_c)$
    \EndIf
\EndFor
\State $v \gets \Verify(\tau_0)$ \Comment{类型检查 / 形式证明检查}
\State $\hat{r} \gets \Render(\tau_0, v)$ \Comment{渲染为自然语言，附 IR 引用}
\State \Return $(v, \hat{r}, \tau_0)$ \Comment{一并返回 IR，供审计}
\end{algorithmic}
\end{algorithm}

算法~\ref{alg:srr} 与朴素「CoT 后比对答案」的差别在于：
它\emph{始终保有一个可检查的形式对象}，并把它作为唯一的真相来源
（single source of truth）。

\subsection{双向一致性}

仅有关注前向映射是不够的：$\Reduce$ 不可逆，使得前向正确性无法被人类确认。
因此我们引入\textbf{双向一致性}（Bidirectional Consistency）约束。

\begin{definition}[双向一致性]\label{def:biconsistency}
设 $\tau = \Reduce(q)$，$\hat{q} = \Render(\tau)$。称 $(q, \tau)$ 满足
双向一致性，当且仅当
\begin{enumerate}[leftmargin=1.8em,itemsep=2pt]
  \item \textbf{可解释性}（interpretability）：$\hat{q}$ 是 $\tau$ 的一个
        忠实视图，即 $\sem{\hat{q}} = \sem{\tau}$；
  \item \textbf{可回溯性}（traceability）：存在从 $\hat{q}$ 的每个语义单元
        到 $\tau$ 的子结构的一一对应（alignment），记作 $\alpha$。
\end{enumerate}
\end{definition}

可回溯性要求给出对齐映射
\begin{equation}\label{eq:align}
  \alpha : \mathrm{Seg}(\hat{q}) \;\longrightarrow\; \mathrm{Sub}(\tau),
\end{equation}
使得人类可以逐句核对「我看到的这句话」对应 IR 的哪一部分。
这直接对抗 \S\ref{sec:bottleneck} 中的第 (2)(3) 条困难：
类型补齐的每一步都必须显式暴露，语义漂移无法隐藏。

\begin{proposition}\label{prop:loss}
若系统满足双向一致性，则任何「类型正确但语义错误」的规约结果
（\S\ref{sec:bottleneck} 第 (2) 条）在原则上可被人类审计发现。
\end{proposition}

\begin{proof}[证明思路]
由定义~\ref{def:biconsistency} 的可回溯性，存在对齐 $\alpha$ 把 $\hat{q}$
的每个语义单元映射到 $\tau$ 的子结构。若 $\tau$ 相对 $q$ 发生了语义漂移，
则必存在某个语义单元（如量词作用域、定义域条件）在 $\hat{q}$ 中
无对应或对应错位；该错位可由审计者沿 $\alpha$ 比对发现。
故漂移不可能在保持 $\hat{q}$ \emph{与原问题逐单元一致}的同时不被察觉。
\end{proof}

\begin{remark}
命题~\ref{prop:loss} 是\emph{可能性}命题而非\emph{复杂性}命题：
它保证漂移\emph{可被发现}，但不保证审计成本可接受。降低审计成本
（例如只审计 $\alpha$ 上置信度最低的若干单元）是一个独立的工程问题。
\end{remark}

% =====================================================================
\section{与现有工作的关系}\label{sec:related}

\subsection{推理范式}

Chain-of-Thought \citep{wei2022cot} 奠定了「显式中间步骤」的范式，
但它把推理与计算混在同一个文本序列中。PoT \citep{chen2023pot} 明确把
\emph{计算}剥离给外部解释器，在多个数据集上平均带来约 $12\%$ 的增益。
Toolformer \citep{schick2023toolformer} 进一步让模型\emph{自监督地}学会
何时调用工具。这三者代表了从「纯文本推理」到「模型 $+$ 外部系统」的
演进，但其外部系统承担的是\emph{计算}，而非\emph{语义规约}。

\subsection{自动形式化与神经定理证明}

autoformalization \citep{wu2022autoformalization} 首次系统性地把
「自然语言 $\to$ 形式规格」作为独立任务，并报告了 $25.3\%$ 的完美翻译率。
DeepSeek-Prover \citep{xin2024deepseekprover} 通过大规模合成数据推进了
形式定理证明的性能。这些工作是本文框架\emph{最直接的先行者}：它们事实上
构建了语义规约层的雏形。本文的差异在于\emph{定位}——我们不把自动形式化
视为一个待优化的翻译模块，而把它视为决定整个系统可靠性的\emph{核心层}，
并据此提出以 Typed IR 为载体的统一设计，以及双向一致性这样的显式契约。

\subsection{神经符号系统}

AlphaGeometry \citep{trinh2024alphageometry} 用神经语言模型与符号演绎引擎
协同求解奥林匹克几何，无需人类演示；FunSearch \citep{romera2024funsearch}
用 LLM 引导的程序搜索在组合数学中发现了新的构造；
AlphaProof \citep{hubert2025alphaproof} 通过强化学习在 Lean 环境中学习证明，
在 2024 IMO 上达到银牌水平；PatternBoost
\citep{charton2024patternboost} 在局部搜索与全局 transformer 之间交替，
为若干长期开放问题找到了更好的构造，包括一个悬置 30 年的猜想的反例。
这些系统是\emph{强证据}：形式化环境 $+$ 可靠验证器 $+$ 搜索/学习，
确实可以产生可验证的新数学。它们的共同特征是\emph{把正确性锚定在形式系统上}——
这与本文的主张一致。本文补充的是：当输入来自自然语言（而非预先形式化的
题面）时，锚定本身需要一整个语义规约层来支撑。

\begin{table}[htbp]
\centering
\caption{与现有工作的定位对比}
\label{tab:compare}
\small
\begin{tabular}{@{}p{3.0cm}p{2.5cm}p{2.6cm}p{2.6cm}p{3.0cm}@{}}
\toprule
\textbf{工作} & \textbf{计算外包} & \textbf{语义规约} & \textbf{形式验证} & \textbf{规约层地位} \\
\midrule
CoT \citep{wei2022cot}            & 否 & 否 & 否 & 无 \\
PoT \citep{chen2023pot}           & 是 & 否 & 否 & 无（仅计算）\\
Toolformer \citep{schick2023toolformer} & 是（多工具）& 否 & 否 & 无 \\
autoformalization \citep{wu2022autoformalization} & 否 & \textbf{是} & 是 & 作为翻译模块 \\
AlphaProof \citep{hubert2025alphaproof} & 是 & 隐含 & \textbf{是} & 隐含在 RL 流程中 \\
\addlinespace
\textbf{本文} & 是（IR 可执行）& \textbf{是（核心）}& 是 & \textbf{独立核心层 $+$ 双向契约} \\
\bottomrule
\end{tabular}
\end{table}

\subsection{本文的差异}

如表~\ref{tab:compare} 所示，本文与已有工作的差别不在「是否使用形式验证」
（多数工作都用了），而在\emph{是否把语义规约确立为独立的、有显式契约的
系统层}。我们把这一层的形式化缺口总结为三点：缺少统一载体（Typed IR）、
缺少双向契约（可解释性 $+$ 可回溯性）、缺少专门的评估指标。

% =====================================================================
\section{应用场景}\label{sec:application}

\subsection{自动定理证明（ATP）}

在证明辅助环境（Lean 4 \citep{demoura2021lean4}、Coq）中，工作流为：
LLM 生成证明草图 $\to$ 规约为 IR $\to$ 展开为形式证明项 $\to$ 内核检查。
与 AlphaProof 式流程的差别在于：我们把「草图 $\to$ IR」这一步显式化，
并允许在人机协作中让人类审计 IR（利用定义~\ref{def:biconsistency} 的
对齐映射）。这对\emph{教学}与\emph{审稿}场景尤其有价值——人类需要理解
证明为什么对，而不只是知道它被接受了。

\subsection{数学解题系统}

对文字应用题（word problems，如 GSM8K \citep{cobbe2021gsm8k}），
标准流程为：题面 $\to$ 类型化程序 $\to$ 执行器求解。
这与 PoT \citep{chen2023pot} 的下游结构相同；本文的补充是在
「题面 $\to$ 程序」之间插入显式类型标注与约束提取（$\mathcal{C}$），
使「单位、定义域、边界条件」这类隐含前提不再被静默丢弃。

\subsection{教学与评估}

在教育场景中，学生解答本身是自然语言。框架可把学生解答规约为 IR，
然后在 IR 层做\emph{结构化}诊断：不只判定对错，而是定位到
「哪一步的类型/约束出了问题」。这比基于最终答案的评分更有信息量，
且由于 IR 可回溯，反馈可以对齐到学生的原句（式~\eqref{eq:align}）。

% =====================================================================
\section{实验设计}\label{sec:experiment}

\subsection{数据集}

选择三个层次分明的数据集，以分离「计算难度」与「规约难度」：
\begin{itemize}[leftmargin=1.6em,itemsep=2pt]
  \item \textbf{GSM8K} \citep{cobbe2021gsm8k}：小学应用题，计算简单，规约中等；
  \item \textbf{MATH} \citep{hendrycks2021math}：竞赛级数学，计算与规约均变难；
  \item \textbf{miniF2F} \citep{zheng2022minif2f} 与 \textbf{PutnamBench}
        \citep{tsoukalas2024putnambench}：直接以形式系统为后端，
        规约难度最高。
\end{itemize}

\subsection{对比设置}

\begin{table}[htbp]
\centering
\caption{实验对比设置}
\label{tab:setup}
\small
\begin{tabular}{@{}p{1.0cm}p{4.6cm}p{7.6cm}@{}}
\toprule
\textbf{组} & \textbf{系统} & \textbf{说明} \\
\midrule
A & 纯 LLM（CoT）& 仅生成层，答案比对式评估 \\
B & LLM $+$ 外部计算（PoT 式） & 生成层 $+$ 计算外包 \\
C & LLM $+$ 形式验证（端到端自动形式化） & 生成层 $+$ 规约 $+$ 验证，但规约层不显式化 \\
\addlinespace
\textbf{D} & \textbf{本文框架（SRR, 算法~\ref{alg:srr}）} & \textbf{显式 Typed IR $+$ 双向一致性契约} \\
\bottomrule
\end{tabular}
\end{table}

\subsection{评估指标}

除常规正确率外，我们提出三个针对规约层的指标：

\paragraph{(1) 正确率（accuracy）.}
$\mathrm{Acc} = \frac{1}{N}\sum_i \mathbb{1}[\hat{a}_i = a_i]$，其中
$\hat{a}_i$ 为系统答案，$a_i$ 为参考答案。这是可比基线。

\paragraph{(2) 可验证率（verifiability）.}
定义
\begin{equation}\label{eq:verifiability}
  \mathrm{Ver} = \frac{\#\{\, i : \Gamma \vdash \tau_i : T_i \ \text{且}\ \mathcal{C}_i \text{ 可满足} \,\}}{N}.
\end{equation}
它度量「有多少比例的输入被成功规约为良类型 IR」。这是该框架下最灵敏的
指标——若 $\mathrm{Ver}$ 低，则系统退化为 A 组，无论生成层多强。

\paragraph{(3) 回译保真度（back-translation fidelity）.}
定义
\begin{equation}\label{eq:fidelity}
  \mathrm{Fid} = \frac{1}{N}\sum_{i}
     \mathrm{Sim}\!\big(q_i,\ \Render(\tau_i)\big),
\end{equation}
其中 $\mathrm{Sim}$ 采用人工评分（盲评）或参考答案约束下的语义等价判定。
$\mathrm{Fid}$ 直接检验双向一致性（定义~\ref{def:biconsistency}）是否成立。

\subsection{假设}

\begin{itemize}[leftmargin=1.6em,itemsep=2pt]
  \item \textbf{H1（归因假设）}：在 MATH 与 miniF2F 上，C 组相对 B 组的
        失败增量主要来自 $\Reduce$（可用 $\mathrm{Ver}$ 的下降解释），
        而非 $\Verify$。
  \item \textbf{H2（框架假设）}：D 组在 $\mathrm{Ver}$ 与 $\mathrm{Fid}$
        上显著优于 C 组，且这一优势在规约难度最高的 PutnamBench 上最大。
  \item \textbf{H3（边界假设）}：在 GSM8K 上，D 组相对 A/B 的优势有限——
        因为简单题目的规约难度低，语义规约层的边际收益小。
\end{itemize}

\begin{remark}
H3 是刻意设置的\emph{证伪性}预测：若框架在所有数据集上都有大幅优势，
反而提示指标设计存在漏洞（例如泄漏了额外信息）。一个诚实的方法论文
应当给出自己\emph{不}适用的边界。
\end{remark}

% =====================================================================
\section{讨论}\label{sec:discussion}

\subsection{局限}

\paragraph{(1) 本框架尚未完成端到端实证.}
本文给出的是设计论证与可检验假设，实验部分（\S\ref{sec:experiment}）是
预注册式方案而非已完成结果。因此，H1--H3 目前处于\emph{待检验}状态。
这是本文最主要的局限。

\paragraph{(2) IR 设计复杂、成本高.}
设计一个既表达力充分又可判定的类型系统，需要领域特定的权衡；
每个新领域（分析、代数、组合）都需要扩展 IR 的构造子。
这与 FunSearch \citep{romera2024funsearch} 式的「轻量 LLM $+$ 现成搜索」
相比，前期投入显著更高。

\paragraph{(3) 语义映射成本高.}
$\Reduce$ 的调用可能涉及多次 LLM 请求 $+$ 形式检查，成本随问题复杂度上升。
在低价值场景（如简单算术）中入不敷出——这正是 H3 所指的边界。

\paragraph{(4) 双向一致性的审计成本.}
命题~\ref{prop:loss} 只保证漂移「可被发现」，不保证廉价。对齐映射 $\alpha$
的人工核查仍需成本。

\subsection{扩展方向}

\paragraph{(1) 通用推理系统.}
把 Typed IR 从数学推广到一般演绎推理（法律、程序验证、科学论证），
其中「命题」与「证据」同样存在自然语言与形式对象之间的鸿沟。

\paragraph{(2) Agent 系统.}
在多智能体架构中，每个 agent 的中间结论都应是可审计的 IR 对象，
而非自由文本。这等价于把本文的双向一致性契约作为 agent 间通信协议，
从而使多 agent 协作\emph{可验证}而非仅仅是「看起来协调」。
这与课程材料中提到的 agent 架构方向一致 \citep{chen2026lectures}。

\paragraph{(3) 与搜索/学习结合.}
IR 的图结构天然适合作为搜索空间：可把 PatternBoost
\citep{charton2024patternboost} 式的局部/全局交替搜索应用到 IR 图上，
用「验证是否通过」作为唯一的奖励信号——因为该信号完全可靠
（由验证层保证），这使得强化学习的目标函数免于被幻觉污染。

% =====================================================================
\section{结论}\label{sec:conclusion}

本文的核心结论可凝练为一个等式：
\begin{equation}\label{eq:conclusion}
  \boxed{\ \text{可靠 AI 推理} \;=\; \text{生成} \;+\; \text{语义规约} \;+\; \text{验证}\ }
\end{equation}
其中三项缺一不可，且\emph{语义规约}是当前最薄弱、也最常被隐式跳过的一环。

本文的关键贡献有两项：
\begin{enumerate}[leftmargin=1.6em,itemsep=2pt]
  \item 明确提出「\textbf{语义规约层}」是可靠 AI 数学的核心瓶颈，
        以三层分解模型给出可检验的归因假设（H1）；
  \item 提供以 \textbf{Typed IR} 为载体的结构化路径，配合
        双向一致性契约与三个新指标，把「可靠性」从口号变为可测量的工程目标。
\end{enumerate}

从更广的视角看，本文主张的转变是：\emph{不要把自然语言当作系统的内部
表示}。自然语言适合作为人机接口，但可靠性必须在结构化的、可验证的
表示上获得。这一原则不限于数学——任何「生成 $\to$ 需要被相信」的场景，
都需要同样的中间层。

% =====================================================================
\section*{致谢}
本文写作大纲参考了课程一手材料《Lectures on AI for Mathematics》
\citep{chen2026lectures}。全部引用均经联网逐条核实；
核实方法与记录见随附的 \code{AI日志.md}（阶段 2）。

% =====================================================================
\bibliography{references}

\end{document}
